\chapter{Predictions, Recovery Benchmarks, and Discriminators}

\noindent\textbf{Status:} \texttt{MODEL-INTERNAL CONSEQUENCES / STANDARD-PHYSICS RECOVERY / DISCRIMINATING TESTS EXPLICITLY SEPARATED}.

This chapter uses the word \emph{prediction} only for a quantity that is fixed before the target observation is inspected. Three classes are kept separate:
\begin{enumerate}
\item \textbf{model-internal consequences}: exact or numerical consequences of the declared IDT model on a stated domain;
\item \textbf{recovery benchmarks}: established physical relations that a viable interface must reproduce;
\item \textbf{discriminating physical tests}: held-out observables for which the IDT construction supplies an independently determined prediction that can differ from an explicit baseline.
\end{enumerate}
Passing a recovery benchmark is necessary compatibility evidence; it is not a new discovery of the recovered relation.

\section{Model-internal and conditional consequences}

\subsection{Positive relational clock ratio}

For any two admitted clocks on the same ordered patch,
\[
\boxed{
N_R(x|r)=\frac{\mathfrak a_x}{\mathfrak a_r}>0.
}
\]
The ratio composes as
\[
N_{x|s}=N_{x|r}N_{r|s}
\]
and is invariant under increasing reparameterisation of the common ordering variable. These are exact consequences of the declared activity measure.

A physical clock claim begins only when the activities can be estimated independently of the target elapsed-time ratio. After a separately established reference calibration,
\[
d\widehat\tau_x=N_R(x|r)\,dt.
\]

\subsection{Relative-information clock potential}

For the admitted local memoryless exponential-rate realization, the standard Kullback--Leibler divergence \cite{KullbackLeibler1951} reduces to
\[
\boxed{
\Phi(N_R)=N_R-1-\ln N_R.
}
\]
It obeys
\[
\Phi(1)=0,
\qquad
\Phi''(N_R)=\frac{1}{N_R^2}>0,
\]
with near-reference expansion
\[
\Phi(1+\epsilon)=\frac12\epsilon^2-\frac13\epsilon^3+O(\epsilon^4).
\]
This is an exact model-class consequence, not yet a physical gravitational prediction.

\subsection{Constructive retrodiction}

On the declared stratified positive-radius decoder domain,
\[
(\alpha,\mathcal W,\rho_1,\ldots,\rho_{N-1},r_N)
\longrightarrow
(r_1,\ldots,r_N)
\longrightarrow
(u_1,\ldots,u_N)
\]
is a constructive inverse chain. The retained radial packet satisfies
\[
\boxed{
N_{\rm radial}=N-1
=\frac12N_{\rm Cartesian}.
}
\]
The associated software tests verify the implementation of this stated map and its rejection conditions. They do not supply independent physical evidence.

\subsection{Chronology-sensitive reference reconstruction}

For E003, a full-rank sensitivity matrix on the admitted reference class permits reconstruction of the hidden kick coordinate, while chronology-destroyed checkpoint controls are required to produce a substantially worse weighted residual
\[
Q=r^T\Sigma_Y^{-1}r.
\]
This is a model-internal observability and reconstruction claim. It is deliberately separated from retrocausal interpretation.

\section{Standard-physics recovery benchmarks}

\subsection{Quantum relative-phase and spectral benchmark}

Once a temporal coordinate is bound to ordinary physical clock time, a stationary discrete quantum system obeys the standard phase-frequency relation
\[
\boxed{
\omega_{mn}=\frac{E_m-E_n}{\hbar}.
}
\]
Recovering this spectrum is a compatibility target, not a novel IDT prediction. The IDT interface fails if its calibrated phase carrier cannot reproduce the independently measured standard relation.

\subsection{Spacetime phase benchmark}

The recombined temporal and spatial branches must recover
\[
\boxed{
\Delta\phi
=
\frac{1}{\hbar}
\left(\mathbf p\cdot\Delta\mathbf x-E\Delta t\right).
}
\]
This Lorentz-covariant phase structure is used as a downstream closure benchmark. Agreement constrains the interface; disagreement falsifies at least one branch or its binding.

\subsection{Einstein and spin-2 benchmark}

The spacetime/source layer must remain compatible with the standard Einstein equation \cite{Einstein1916} and with the linearized transverse-traceless response of a conserved directional source. In particular,
\[
h_{ij}^{TT}
=
\frac{4G}{c^4r}\mathcal T_{ij}^{TT}
\]
is treated as a recovery relation. The novel content, if any, must reside in an independently specified IDT source predictor rather than in this standard propagation equation.

\subsection{Horizon thermodynamic benchmark}

For a non-extremal Euclidean horizon,
\[
\beta_H=\frac{2\pi}{\kappa_H},
\qquad
\boxed{
T_H=\frac{\hbar\kappa_H}{2\pi c k_B}.
}
\]
These Hawking/Gibbons--Hawking relations are standard \cite{Hawking1975,GibbonsHawking1977}. Integer-spin periodicity and half-integer-spin antiperiodicity are likewise boundary-condition recovery targets. Their reproduction is a consistency requirement, not a new IDT prediction.

\section{Discriminating physical test D1: independent activity-to-clock prediction}

Let \(Z_x,Z_r\) denote independently measured relational observables that do not include the target clock-rate result. Before the held-out clock comparison is inspected, freeze
\[
\boxed{
\widehat N_R^{\,\rm IDT}
=
f_{\theta}(Z_x,Z_r),
}
\]
including the parameter vector \(\theta\), preprocessing, exclusions, and uncertainty model. The target measurement is
\[
R_{\rm clock}
=
\frac{d\tau_x}{d\tau_r}.
\]
The primary residual is
\[
\boxed{
\Delta_N
=
R_{\rm clock}
-
\widehat N_R^{\,\rm IDT}.
}
\]

The preregistered baseline \(N_R^{(0)}\) may be the standard geometric prediction or another declared physical model. Discrimination requires held-out comparison of the predictive residuals or likelihoods,
\[
\mathcal L_{\rm IDT}(R_{\rm clock})
\quad\text{versus}\quad
\mathcal L_{0}(R_{\rm clock}),
\]
with all shared nuisance parameters fitted only on a disjoint calibration set.

The test is \emph{not} discriminating if \(R_{\rm clock}\), \(N_R^{(0)}\), or an algebraically equivalent target quantity is used to define \(Z\), \(\theta\), or \(\mathfrak a_x/\mathfrak a_r\). In that case the result is classified as calibration/consistency. If the independently generated IDT predictor is mathematically identical to the baseline on the tested domain, the outcome is a recovery result rather than evidence favoring either description.

\section{Discriminating physical test D2: half-interface \(4\pi\) structure}

The IDT half-interface kernel is
\[
\boxed{
\mathcal D_{1/2}(p,\Delta\tau)
=
1+2\sqrt{p(1-p)}\cos(\Delta\tau/2).
}
\]
At equal weights it has the exact model null
\[
p=\frac12,
\qquad
\Delta\tau\equiv2\pi\pmod{4\pi}.
\]

A physical test must freeze the map from laboratory control variables to \(\Delta\tau\) before the target scan. The confirmatory acquisition spans a full \(4\pi\) interval and includes:
\begin{enumerate}
\item an equal-weight condition \(p=1/2\);
\item weight-detuned controls \(p\neq1/2\);
\item phase-detuned controls around \(2\pi\);
\item the appropriate conventional-interference baseline expressed in the same laboratory control coordinate;
\item blinded analysis of the null location and periodicity.
\end{enumerate}

A \(4\pi\) pattern is discriminating only if \(\Delta\tau\) has an independently established physical calibration. Redefining an ordinary \(2\pi\) phase by an arbitrary factor of two would merely relabel the coordinate and is explicitly excluded as evidence.

\section{Discriminating physical test D3: future-conditioned committed record}

The prospective E004 protocol targets
\[
\boxed{
P(X_{\mathrm{past}}\mid Y_{\mathrm{future}})
\stackrel{?}{=}
P(X_{\mathrm{past}})
}
\]
as the null hypothesis. The alternative is a dependence that remains after the future condition is generated strictly after the earlier record has been irreversibly committed.

The preregistration in Chapter~9 fixes the temporal order, valid-trial target, stopping rule, frozen score, permutation analysis, channel audit, and control datasets before confirmatory acquisition. A candidate effect is not promoted by statistical significance alone. It must survive:
\begin{enumerate}
\item commitment and chronology verification;
\item classical information-channel and shared-state audits;
\item sham, leakage-positive, and scheduler/shared-clock controls;
\item frozen exclusions and analysis rules;
\item independent replication using the same causal-order contract.
\end{enumerate}
This is the strongest causal discriminator in the current programme because its null and chronology are defined without using the observed effect.

\section{Prediction and promotion discipline}

A physical prediction is considered genuinely discriminating only when all four conditions hold:
\begin{enumerate}
\item \textbf{independence}: the predictor is fixed without access to the target outcome;
\item \textbf{contrast}: an explicit competing/null model makes a distinct prediction on the same observable;
\item \textbf{uncertainty}: calibration, nuisance parameters, and measurement uncertainty are propagated into the comparison;
\item \textbf{replicability}: the analysis contract can be rerun on held-out or independently acquired data.
\end{enumerate}

Exact identities and recovery benchmarks remain scientifically useful even when they are not discriminators. This classification prevents compatibility with established physics, successful software tests, and genuinely new empirical content from being counted as the same kind of evidence.
